\documentclass[../main.tex]{subfiles}

\begin{document}
\setcitestyle{square}
\maketitle
\thispagestyle{empty}
\rule{411.8pt}{1pt}
\begin{abstract}
\noindent The NA62 experiment is a fixed target experiment located in the north area of CERN, which aims to obtain precision measurements of charged kaon decays. Its main aim is to measure the branching ratio of the ultra-rare decay $K^+ \rightarrow \pi^+ \nu \bar{\nu}$ to a 10\% uncertainty to match theoretical results. Results from run 1 data taken between 2016-2018 have been analysed, and 20 candidates have been identified, resulting in a branching ratio $\mathcal{B}(K^+ \rightarrow \pi^+ \nu \bar{\nu}) = (11.0^{+4.0}_{-3.5} \pm 0.3 ) \times 10^{-11}$. 
\newline The current NA62 experiment is in its second year of run 2 data taking and has been approved up to CERN's Long Shutdown 3 in 2025. After this, there is a program for High-Intensity kaon Experiments (HIKE) which will involve both charged and neutral kaon experiments. This report covers the improvements to the NA62 offline software to aid the development of new detector setups. Using this software, a symmetric NA62 setup has been proposed to allow for a beam to travel down the centre of the sub-detectors, allowing for both $K^+$ and $K_L$ beams. The full Software package has allowed for momentum and squared missing mass resolutions to be calculated to evaluate the new detector setup
\newline\linebreak
\newline\textbf{Student ID:} 1777911
\newline\textbf{Key Words:} 
\end{abstract}



\rule{411.8pt}{1pt}
$^1$jack.sanders@cern.ch
\end{document}